\scanpage{74}

% L01 -- top
Reformulation des données lorsque $c\notin C$ en \unsure{tte fibre|ts fibres} i.e. $\rho$ inv ($n$ impair, $\rho'$ inv, ou $n$ pair, $\rho'$ inv. et 2 inv.)

% L02 -- upper third, A
\[
\text{A)}\quad \varphi_G:G\longrightarrow\mathrm{Aut}(X)\quad
\text{ou encore}\quad(\tau_i^X)_{0\le i\le n}
\]

% L03 -- upper third
est \unsure{remplacé} par

% L04 -- upper third, second action
\[
\varphi_G:G\longrightarrow\mathrm{Aut}(\unsure{W|V})\quad
\text{ou encore}\quad(\tau_i^V)_{0\le i\le n}
\]

% L05 -- above numbered list
Conditions

% L06 -- first condition
$1^\circ)$ les $\tau_i^V$ pseudo-réflexions (\unsure{d'où|définit|c.à.d.} $L_i\subset V_{\unsure{i}},\ H_i\subset V_i$)

% L07 -- left of second condition, diagonal
% slanted
\margin{\deleted{\unsure{déf. des fais}}}

% L08 -- second condition
\[
2^\circ)\quad V=\bigoplus_{0\le i\le n}L_i,\quad V^\vee\simeq
\bigoplus\underbrace{H_i^0}_{L'_0}\quad\text{i.e.}\quad
\bigcap_{0\le\unsure{j|i}\le n}H_i=0\quad\text{chaque fibre}
\]

% L09 -- third condition
\[
3^\circ)\quad L_i\subset H_{i+1}\quad\text{chaque fibre}\qquad0\le i\le n-1
\]

% L10 -- cancelled fourth condition
\deleted{\[
4^\circ)\quad\text{\unsure{Posant}}\quad D_0=\bigcap_{1\le i\le n}H_i,\quad\text{\unsure{il est|est}}
\]}

% L11 -- middle and lower part, B
B) La donnée de $s_0\in\unsure{E\subset X|\Gamma(X)}$ \deleted{\unsure{en}} \unsure{donnant lieu avec|donnant lieu à} $\varphi_G$ à une \unsure{réal. géom} non dég, est \deleted{\unsure{rempl}} exprimée par

% L12 -- middle, base
\[
\text{\uline{$\widetilde{s}_0$ base de}}\quad D_0=\bigcap_{1\le i\le n}H_i
\]

% L13 -- lower third
qui épingle la situation. \deleted{On obtient} [Il revient au $=$ de se donner une base \deleted{C) Les formes quadratiques} $a'_0$ de $L'_0$ (car $D_0\otimes L'_0\xrightarrow{\sim}\mathcal O_S$, par

% L14 -- lower third, boxed
\[
\boxed{\langle\widetilde{s}_0,a'_0\rangle=1}
\]

% L15 -- foot
ce qui fixe bien $a_0$ de $L_0$ par $\tau_0=\mathrm{id}_V+a'_0\otimes a_0$, et de proche en proche on détermine les $a'_i,a_i$ par $\langle a_i,a'_{i+1}\rangle=1,\qquad\tau_i=\mathrm{id}+a'_i\otimes a_i$ ]
% archive sheet number: 73
